\section{组织平权：为什么 Agent 研究需要多样性}

访谈中“组织平权”不是政治口号，而是研究组织方法。罗福莉认为，做后训练需要 diversity，让预训练的人参与后训练是很好的补充。Agent 系统涉及模型、数据、RL、产品、工程、成本、安全和用户环境，单一背景团队很难覆盖全部问题。

\begin{figure}[H]
\centering
\includegraphics[width=0.92\textwidth]{org-equality.png}
\caption{Agent 研究中的组织平权：多职能进入同一研究闭环。}
\end{figure}

\begin{knowledgebox}{读图：组织平权不是平均主义}
图中 pre-train、post-train/RL、product/UX、infra/systems、leadership 都连向 Agent research。它说明不同职能都需要进入研究闭环，而不是由某个单一小组垄断判断。平权指信息和责任进入同一系统，而不是每个人做同样工作。
\end{knowledgebox}

\subsection{群体智能如何提升 Agent 框架}

春节期间，罗福莉推动团队集体使用 OpenClaw。她强调，一个人的想象力有限；当团队成员在大群里展示“原来还能这么用”时，会激发其他人的想象力。群体智能不是抽象概念，而是通过共享使用经验、互相刺激场景想象、快速发现框架能力和短板来提升研究速度。

\begin{importantbox}{组织也是训练环境}
如果 Agent 能从环境中学习，团队也能从共同使用环境中学习。让团队共同进入新工具、新范式、新任务环境，本质上是在给组织做一次 post-train。
\end{importantbox}

\subsection{文化与价值观}

视频描述强调，这次技术巨变中 AI Lab 的根基是文化与价值观。罗福莉后面也谈到，自己每天做的事要让世界变好一点，让无聊工作被替代，让人有时间做更有价值的事。这个价值观会影响团队如何选择任务、如何定义 benchmark、如何配置资源。

\begin{warningbox}{价值观不是装饰}
当模型越来越能替代人类工作时，团队选择优化什么、替代什么、保留什么，会变成技术决策的一部分。没有价值观约束，Agent 能力越强，风险也可能越大。
\end{warningbox}

\subsection{本章小结}

Agent 时代的组织能力不只是招更多人，而是让不同背景的人共同进入研究闭环。组织平权、群体智能和价值观，决定一个实验室能否快速拥抱新范式。

